% Generated draft: requires mathematical and source review.
\documentclass[11pt]{article}
\usepackage[margin=25mm]{geometry}
\usepackage{fontspec}
\setmainfont{Latin Modern Roman}
\usepackage{amsmath,amssymb,mathtools,unicode-math}
\setmathfont{Latin Modern Math}
\usepackage{enumitem,needspace,xurl,hyperref}
\setlength{\parindent}{0pt}
\setlength{\parskip}{5pt}
\setlength{\emergencystretch}{3em}
\newcommand{\N}{\mathbb{N}}
\newcommand{\Z}{\mathbb{Z}}
\newcommand{\Q}{\mathbb{Q}}
\newcommand{\R}{\mathbb{R}}
\newcommand{\C}{\mathbb{C}}
\newcommand{\F}{\mathbb{F}}
\newcommand{\A}{\mathbb{A}}
\newcommand{\PP}{\mathbb{P}}
\newcommand{\E}{\mathbb{E}}
\newcommand{\eps}{\varepsilon}
\newcommand{\dd}{\,\mathrm{d}}
\begin{document}
\section*{UnsolvedMath ID: 30001048}
\textbf{Source problem code:} catalogue entry
\subsection{Definitions and mathematical statement}
Let $C$ be the compression body from the source's hyperbolic tunnel-number-one manifold after removal of its handlebody side. It has negative boundary $\partial_-C\simeq T^2$ and positive boundary $\partial_+C$ of genus two; topologically this is the type obtained from $T^2\times[0,1]$ by attaching one $1$-handle along the upper boundary. Let $\tau$ be the source's core tunnel arc through that compressible handle. The exact proper-arc endpoints and product extensions of the handle core are not described in the extract. The conjecture asks whether \[\forall\text{ geometrically finite hyperbolic structures }h\text{ on }C:\qquad\tau\text{ is isotopic to an embedded }h\text{-geodesic arc}.\] The standard geometric interpretation is a marked complete hyperbolic structure on $\operatorname{int}C$, with the torus end treated as a cusp and the fundamental group acting discretely and faithfully on $\mathbb H^3$ with geometrically finite quotient. The source does not specify the parabolic/ending conventions, including whether additional accidental parabolics are allowed. If the core tunnel has its ends on the torus cusp, its geodesic is understood as a cusp-to-cusp ideal geodesic, with proper isotopy or a compatible truncation convention. If other endpoints are intended, a different relative-endpoint formulation is needed; this is not resolved by the extract. The assertion is isotopy to a geodesic, retaining the embedded arc's topological class, rather than only existence of a geodesic representative in a homotopy class. The given tunnel-number-one origin and every-structure quantifier are not replaced by a statement about a special family of structures. The supplied 2026-08-21 triage labels the question open, says the report context identifies the core-tunnel conjecture despite an extracted heading, and reports proofs for broad families without the general assertion.
\subsection{Short English statement}
Must the core tunnel of the specified torus-to-genus-two compression body be isotopic to a geodesic in every geometrically finite hyperbolic structure? The arc endpoints, cusp conventions, and relative isotopy formulation must be recovered from the report.
\subsection{Sources}
\begin{itemize}\item Local UnsolvedMath registry, permanent ID 30001048, source problem catalogue entry, \emph{Geodesicity of Core Tunnels in Compression Bodies}, proposed year 2008; no individual proposer supplied.\item Supplied citation: \emph{Topologie} (2009), Oberwolfach Reports, DOI 10.4171/owr/2008/43, pp.\ 2419--2476: \url{https://doi.org/10.4171/owr/2008/43}. Supplied literature label: original report.\item Supplied literature label: Lackenby--Purcell, \url{https://arxiv.org/abs/1302.3652}.\item Literature assessment date: 2026-08-21. The status and progress claims are supplied dated curation, not independently verified here.\end{itemize}
\end{document}
